\chapter*{Abstract}

This project develops a numerical method for identifying multiple simultaneous leaks in a single pipeline, where leak identification refers to both localization and size estimation. The method relies solely on steady-state inlet and outlet head and flow measurements to support software-based leak identification in water distribution systems with limited monitoring equipment. Pipeline leakage causes substantial water loss, increased operational costs, and infrastructure risk, making practical methods for leak identification important. However, existing steady-state methods are largely limited to one or two leaks, while methods for several leaks typically require transient data or additional sensors.\\

To address this gap, the leak identification problem was formulated as a nonlinear residual problem and solved using a Levenberg--Marquardt method with an analytically derived Jacobian. The method was evaluated through an extensive simulated test framework and on published simulated two-leak experiments. The results show that the method performs reliably for two and three leaks and can still provide useful localization for four leaks in many cases. A condition number analysis further showed that the inverse problem becomes rapidly more ill-conditioned as the number of leaks increases. Combined with the empirical results, this suggests that identifying five or more simultaneous leaks is likely close to the upper limit of what is numerically feasible using only steady-state endpoint measurements.\\

These findings show that steady-state endpoint measurements can identify more simultaneous leaks than comparable published steady-state methods. At the same time, they reveal a practical limit imposed by the intrinsic conditioning of the inverse problem.